\documentclass[aspectratio=169]{beamer}
% Shared Beamer setup for Fall 2026 course slides.
\mode<presentation>{
  \usetheme{Madrid}
  \usecolortheme{default}
}

\definecolor{CalPolyGreen}{HTML}{154734}
\definecolor{CalPolyGold}{HTML}{BD8B13}
\definecolor{DeckInk}{HTML}{1F2933}
\definecolor{DeckMuted}{HTML}{5F6B7A}

\setbeamercolor{structure}{fg=CalPolyGreen}
\setbeamercolor{palette primary}{bg=CalPolyGreen,fg=white}
\setbeamercolor{palette secondary}{bg=DeckInk,fg=white}
\setbeamercolor{palette tertiary}{bg=CalPolyGold,fg=DeckInk}
\setbeamercolor{frametitle}{bg=CalPolyGreen,fg=white}
\setbeamercolor{title}{fg=CalPolyGreen}
\setbeamercolor{normal text}{fg=DeckInk,bg=white}
\setbeamercolor{block title}{bg=CalPolyGreen,fg=white}
\setbeamercolor{block body}{bg=black!3,fg=DeckInk}

\setbeamertemplate{navigation symbols}{}
\setbeamertemplate{footline}[frame number]
\setbeamertemplate{itemize item}{\color{CalPolyGold}$\blacktriangleright$}
\setbeamertemplate{itemize subitem}{\color{CalPolyGreen}$\bullet$}

\newcommand{\CourseNumber}{Course}
\newcommand{\CourseTitle}{Course Title}
\newcommand{\LectureNumber}{0}
\newcommand{\LectureTitle}{Lecture Title}
\newcommand{\LectureDate}{Date}
\newcommand{\CourseUnit}{Unit}

\newcommand{\coursenumber}[1]{\renewcommand{\CourseNumber}{#1}}
\newcommand{\coursetitle}[1]{\renewcommand{\CourseTitle}{#1}}
\newcommand{\lecturenumber}[1]{\renewcommand{\LectureNumber}{#1}}
\newcommand{\lecturetitle}[1]{\renewcommand{\LectureTitle}{#1}}
\newcommand{\lecturedate}[1]{\renewcommand{\LectureDate}{#1}}
\newcommand{\courseunit}[1]{\renewcommand{\CourseUnit}{#1}}

\author{Kirk Duran}
\institute{California Polytechnic State University, San Luis Obispo}
\date{\LectureDate}

\newcommand{\makedecktitle}{
  \begin{frame}[plain]
    \vspace{0.25in}
    {\usebeamerfont{title}\color{CalPolyGreen}\LectureTitle\par}
    \vspace{0.18in}
    {\large \CourseNumber{}: \CourseTitle\par}
    \vspace{0.08in}
    {\large Lecture \LectureNumber{} -- \LectureDate\par}
    \vfill
    {\small Kirk Duran \textbar{} Cal Poly, San Luis Obispo\par}
  \end{frame}
}

\newcommand{\placeholdernote}{
  \begin{block}{Placeholder}
    Replace these bullets with the final lecture material, examples, and in-class activities.
  \end{block}
}

\AtBeginSection[]{
  \begin{frame}{Roadmap}
    \tableofcontents[currentsection]
  \end{frame}
}

\graphicspath{{../assets/lecture-01/}}

\coursenumber{CS 4880}
\coursetitle{Artificial Intelligence}
\lecturenumber{1}
\lecturetitle{What Is Artificial Intelligence?}
\lecturedate{August 24, 2026}
\courseunit{0. Introduction and Framing}

\begin{document}

\makedecktitle

% ============================================================
% Opening
% ============================================================

\begin{frame}{Artificial Intelligence}
  \begin{keyidea}
    AI has more public attention than ever. Before we study the techniques, we need to decide what we actually mean by ``artificial intelligence.''
  \end{keyidea}
  \vspace{0.06in}
  \begin{itemize}
    \item AI is no longer confined to science fiction or research laboratories.
    \item It increasingly shapes work, education, medicine, media, science, and everyday decisions.
    \item The term ``AI'' now refers to many different systems with very different capabilities.
  \end{itemize}
  \vspace{0.06in}
  \begin{activity}
    When you hear ``AI,'' what system or capability comes to mind first?
  \end{activity}
\end{frame}

\begin{frame}{How We Will Talk In This Course}
  \begin{columns}[T,onlytextwidth]
    \begin{column}{0.48\textwidth}
      \begin{block}{Socratic Method}
        We will often begin with a question before we introduce the formal answer.
      \end{block}
      \begin{itemize}
        \item I will ask for your intuition.
        \item I may call on people around the room.
        \item ``I do not know'' is a valid starting point.
      \end{itemize}
    \end{column}
    \begin{column}{0.46\textwidth}
      \begin{keyidea}
        The goal is not to guess what I am thinking. The goal is to expose assumptions, compare explanations, and make our reasoning precise.
      \end{keyidea}
    \end{column}
  \end{columns}
\end{frame}

\begin{frame}{Today}
  \begin{itemize}
    \item What do we mean by artificial intelligence?
    \item What do we mean by intelligence?
    \item Where did AI come from?
    \item What kinds of methods count as AI?
    \item What questions will this course try to answer?
    \item What larger technical and social questions surround the field?
  \end{itemize}
\end{frame}

% ============================================================
\sectionframe{Part 1: What Is Artificial Intelligence?}
% ============================================================

\begin{frame}{A Working Definition}
  \begin{center}
    {\Huge What should count as AI?}
  \end{center}
  \vspace{0.18in}
  \begin{activity}
    Propose a one-sentence definition. What must a system be able to do before you would call it intelligent?
  \end{activity}
\end{frame}

\begin{frame}[shrink=4]{Definitions Emphasize Different Things}
  \begin{columns}[T,onlytextwidth]
    \begin{column}{0.31\textwidth}
      \begin{block}{AIMA}
        Focuses on \textbf{agents}: systems that perceive an environment and choose actions.
      \end{block}
    \end{column}
    \begin{column}{0.31\textwidth}
      \begin{block}{Industry}
        Often emphasizes machines performing tasks associated with human reasoning, prediction, or decision-making.
      \end{block}
    \end{column}
    \begin{column}{0.31\textwidth}
      \begin{block}{Engineering}
        Often frames AI as building systems that exhibit useful intelligent behavior.
      \end{block}
    \end{column}
  \end{columns}
  \vspace{0.15in}
  \begin{keyidea}
    There is no single universally accepted definition. Different definitions emphasize behavior, mechanisms, goals, or comparison with humans.
  \end{keyidea}
\end{frame}

\begin{frame}[shrink=3]{More Ways To Define AI}
  \begin{columns}[T,onlytextwidth]
    \begin{column}{0.31\textwidth}
      \begin{block}{Human-like Tasks}
        Systems that perform tasks normally associated with human intelligence: perception, language, reasoning, and decision-making.
      \end{block}
    \end{column}
    \begin{column}{0.31\textwidth}
      \begin{block}{Goal-Directed Systems}
        Machines that make predictions, recommendations, or decisions in pursuit of human-defined objectives.
      \end{block}
    \end{column}
    \begin{column}{0.31\textwidth}
      \begin{block}{Learning And Problem Solving}
        Computational systems that learn, reason, search, plan, or otherwise solve problems.
      \end{block}
    \end{column}
  \end{columns}
\end{frame}

\begin{frame}{Our Working Definition}
  \begin{cpdefinition}
    Artificial intelligence is the study and engineering of computational systems that perceive, reason, learn, decide, or act in ways that accomplish goals we associate with intelligent behavior.
  \end{cpdefinition}
  \vspace{0.14in}
  \begin{itemize}
    \item The system does not need to think exactly like a human.
    \item The system does not need to be generally intelligent.
    \item Intelligence can appear in narrow, well-defined tasks.
    \item We care about both \textbf{what} the system can do and \textbf{how} it does it.
  \end{itemize}
\end{frame}

\begin{frame}{Why Are You Studying AI?}
  \begin{center}
    {\Huge I want to hear from you.}
  \end{center}
  \vspace{0.16in}
  \begin{activity}
    What made you enroll in this course? Curiosity, career goals, research, fear of missing out, wanting to understand modern systems, something else?
  \end{activity}
\end{frame}

\begin{frame}[shrink=4]{Why I Am Interested In AI}
  \begin{columns}[T,onlytextwidth]
    \begin{column}{0.31\textwidth}
      \begin{block}{Automation}
        If a machine can reliably do repetitive work, why should a human have to do it?
      \end{block}
    \end{column}
    \begin{column}{0.31\textwidth}
      \begin{block}{Optimism}
        AI can amplify human capability and potentially improve health, science, accessibility, education, and productivity.
      \end{block}
    \end{column}
    \begin{column}{0.31\textwidth}
      \begin{block}{Philosophy}
        Building artificial minds forces us to ask what intelligence is, what understanding is, and what is uniquely human.
      \end{block}
    \end{column}
  \end{columns}
\end{frame}

% ============================================================
\sectionframe{Part 2: Where Did AI Come From?}
% ============================================================

\begin{frame}{A Very Short History}
  \begin{keyidea}
    AI has repeatedly moved through cycles of ambitious goals, technical breakthroughs, bottlenecks, disappointment, and renewed progress.
  \end{keyidea}
  \vspace{0.08in}
  \begin{itemize}
    \item \textbf{1950s--1960s:} symbolic reasoning, search, theorem proving.
    \item \textbf{1970s--1980s:} expert systems and knowledge engineering.
    \item \textbf{1990s--2000s:} probabilistic reasoning, data-driven learning, stronger computing infrastructure.
    \item \textbf{2010s--present:} deep learning, large-scale models, reinforcement learning, and generative AI.
  \end{itemize}
\end{frame}

\begin{frame}{The Recurring Question}
  \begin{center}
    {\Huge What parts of intelligent behavior can be represented computationally?}
  \end{center}
  \vspace{0.16in}
  \begin{keyidea}
    The tools change, but AI keeps returning to this question.
  \end{keyidea}
\end{frame}

\begin{frame}{20th-Century Beginnings}
  \begin{columns}[T,onlytextwidth]
    \begin{column}{0.47\textwidth}
      \begin{block}{Bold Ambition}
        Build machines that could reason and act in ways associated with human intelligence.
      \end{block}
      \begin{block}{Scientific Motivation}
        Formalize reasoning, problem solving, language, learning, and decision-making.
      \end{block}
    \end{column}
    \begin{column}{0.47\textwidth}
      \begin{block}{Engineering Motivation}
        Automate difficult tasks and extend what machines can do.
      \end{block}
      \begin{block}{Social Context}
        War, economic competition, research funding, and technological optimism influenced the field's priorities.
      \end{block}
    \end{column}
  \end{columns}
\end{frame}

\begin{frame}{1950: Turing Changes The Question}
  \begin{columns}[T,onlytextwidth]
    \begin{column}{0.58\textwidth}
      \begin{itemize}
        \item Alan Turing asked whether machines could exhibit intelligent behavior.
        \item He proposed an operational test based on text conversation.
        \item The focus shifts from defining an invisible mental property to examining observable behavior.
      \end{itemize}
      \vspace{0.10in}
      \begin{keyidea}
        When a concept is hard to define directly, ask what evidence would demonstrate it.
      \end{keyidea}
    \end{column}
    \begin{column}{0.35\textwidth}
      \begin{block}{Question}
        If a machine can convincingly imitate intelligent conversation, is that enough to call it intelligent?
      \end{block}
    \end{column}
  \end{columns}
\end{frame}

\begin{frame}{1956: Artificial Intelligence Becomes A Field}
  \begin{itemize}
    \item John McCarthy helped organize the Dartmouth summer research project where the term \textbf{artificial intelligence} was introduced.
    \item Researchers proposed that aspects of learning and intelligence might be described precisely enough for machines to simulate them.
    \item The field began with unusually ambitious expectations about general-purpose machine intelligence.
  \end{itemize}
  \vspace{0.14in}
  \begin{activity}
    Is it easier to build a machine that is excellent at one task, or a machine that is broadly competent across many tasks? Why?
  \end{activity}
\end{frame}

% ============================================================
\sectionframe{Part 3: What Do We Mean By Intelligence?}
% ============================================================

\begin{frame}{Intelligence}
  \begin{center}
    {\Huge What does intelligence mean?}
  \end{center}
  \vspace{0.18in}
  \begin{activity}
    Name one capability you think is essential to intelligence. Could a system possess that capability without possessing the others?
  \end{activity}
\end{frame}

\begin{frame}[shrink=4]{Behavior Is Not The Same As Understanding}
  \begin{columns}[T,onlytextwidth]
    \begin{column}{0.5\textwidth}
      \begin{block}{The Chinese Room}
        John Searle's thought experiment asks us to imagine a person manipulating unfamiliar symbols by following formal rules.
      \end{block}
      \begin{itemize}
        \item Inputs arrive as symbols.
        \item Rules determine which symbols to return.
        \item To an outsider, the responses may appear meaningful.
      \end{itemize}
    \end{column}
    \begin{column}{0.44\textwidth}
      \begin{keyidea}
        The philosophical challenge: can correct symbol manipulation produce genuine understanding, or only the appearance of understanding?
      \end{keyidea}
      \vspace{0.10in}
      \begin{activity}
        Does this argument become more or less convincing when applied to modern language models?
      \end{activity}
    \end{column}
  \end{columns}
\end{frame}

\begin{frame}{AI Is More Than ChatGPT}
  \begin{keyidea}
    Large language models are one part of AI. The field includes systems with very different goals, representations, and mechanisms.
  \end{keyidea}
  \vspace{0.12in}
  \begin{columns}[T,onlytextwidth]
    \begin{column}{0.30\textwidth}
      \begin{block}{Simple Automation}
        Sensors, rules, controllers, and reactive systems.
      \end{block}
    \end{column}
    \begin{column}{0.30\textwidth}
      \begin{block}{Modern AI}
        Search, planning, probabilistic models, machine learning, neural networks, robotics, and language systems.
      \end{block}
    \end{column}
    \begin{column}{0.30\textwidth}
      \begin{block}{General Intelligence}
        Hypothetical systems with broad, flexible competence across many domains.
      \end{block}
    \end{column}
  \end{columns}
\end{frame}

\begin{frame}[shrink=5]{AI Is A Toolbox Of Mechanisms}
  \begin{columns}[T,onlytextwidth]
    \begin{column}{0.18\textwidth}
      \begin{block}{Search}
        Explore possible states or actions.
      \end{block}
    \end{column}
    \begin{column}{0.18\textwidth}
      \begin{block}{Logic}
        Represent facts and derive conclusions.
      \end{block}
    \end{column}
    \begin{column}{0.18\textwidth}
      \begin{block}{Probability}
        Reason under uncertainty.
      \end{block}
    \end{column}
    \begin{column}{0.18\textwidth}
      \begin{block}{Learning}
        Improve behavior from data or experience.
      \end{block}
    \end{column}
    \begin{column}{0.18\textwidth}
      \begin{block}{Hybrid Systems}
        Combine multiple representations and methods.
      \end{block}
    \end{column}
  \end{columns}
  \vspace{0.16in}
  \begin{keyidea}
    AI is not one algorithm. It is a collection of approaches for building systems that behave intelligently.
  \end{keyidea}
\end{frame}

% ============================================================
\sectionframe{Part 4: What Will We Study?}
% ============================================================

\begin{frame}[shrink=4]{AI Course Roadmap}
  \begin{columns}[T,onlytextwidth]
    \begin{column}{0.31\textwidth}
      \begin{block}{Unit 1: Intelligent Agents}
        Perception, actions, environments, rational behavior.
      \end{block}
      \begin{block}{Unit 2: Search}
        State spaces, heuristics, path finding, optimization.
      \end{block}
      \begin{block}{Unit 3: Adversarial Games}
        Minimax, alpha--beta pruning, Monte Carlo search.
      \end{block}
    \end{column}
    \begin{column}{0.31\textwidth}
      \begin{block}{Unit 4: Knowledge-Based Systems}
        Logic, representation, inference, planning.
      \end{block}
      \begin{block}{Unit 5: Uncertain Knowledge}
        Probability, Bayes, graphical models, tracking.
      \end{block}
    \end{column}
    \begin{column}{0.31\textwidth}
      \begin{block}{Unit 6: Learning}
        Machine learning, neural networks, reinforcement learning.
      \end{block}
      \begin{block}{Unit 7: Language-Based Systems}
        Language models and modern generative AI.
      \end{block}
    \end{column}
  \end{columns}
\end{frame}

\begin{frame}[shrink=3]{Questions We Hope To Answer}
  \begin{columns}[T,onlytextwidth]
    \begin{column}{0.48\textwidth}
      \begin{itemize}
        \item How can an agent perceive its surroundings and choose useful actions?
        \item How can a machine explore alternatives through trial and error?
        \item How can a program search large spaces efficiently?
        \item How can computers represent knowledge and draw conclusions?
      \end{itemize}
    \end{column}
    \begin{column}{0.48\textwidth}
      \begin{itemize}
        \item How can AI reason when information is uncertain?
        \item What allows a system to learn from data or experience?
        \item How do modern language and generative systems work at a high level?
        \item Which problems remain fundamentally difficult?
      \end{itemize}
    \end{column}
  \end{columns}
\end{frame}

\begin{frame}{The Ambitious Goal}
  \begin{center}
    {\Huge How general can artificial intelligence become?}
  \end{center}
\end{frame}

\begin{frame}{Artificial General Intelligence}
  \begin{cpdefinition}
    Artificial General Intelligence (AGI) usually refers to a hypothetical system with broad, flexible competence across many intellectual tasks rather than exceptional performance on only one narrow problem.
  \end{cpdefinition}
  \vspace{0.14in}
  \begin{itemize}
    \item Modern AI is impressive but uneven: systems can be extraordinary in one setting and brittle in another.
    \item AGI is a research aspiration and a contested concept, not a capability we should assume has already been achieved.
    \item Studying narrow AI still teaches us the mechanisms from which more general systems may be built.
  \end{itemize}
\end{frame}

% ============================================================
\sectionframe{Part 5: What Should We Want From AI?}
% ============================================================

\begin{frame}{Is This A Good Thing Or A Bad Thing?}
  \begin{center}
    {\Huge I want your opinion.}
  \end{center}
  \vspace{0.16in}
  \begin{activity}
    If AI becomes dramatically more capable over your lifetime, what is one outcome you hope for and one outcome you worry about?
  \end{activity}
\end{frame}

\begin{frame}[shrink=4]{Automation, Power, And Control}
  \begin{columns}[T,onlytextwidth]
    \begin{column}{0.48\textwidth}
      \begin{block}{The Promise}
        Automation can reduce repetitive labor, improve access to expertise, accelerate discovery, and expand human capability.
      \end{block}
      \begin{block}{The Risk}
        The same systems can concentrate decision-making power, enable surveillance, encode bias, or make institutions harder to challenge.
      \end{block}
    \end{column}
    \begin{column}{0.46\textwidth}
      \begin{exampleblockcp}
        A useful political-philosophy lens: systems can trade autonomy for coordination, efficiency, or security.
      \end{exampleblockcp}
      \vspace{0.10in}
      \begin{activity}
        Who should be allowed to delegate decisions to AI, and which decisions should remain meaningfully contestable by humans?
      \end{activity}
    \end{column}
  \end{columns}
\end{frame}

\begin{frame}{Case Study Prompt: Replacing Human Work}
  \begin{keyidea}
    Automating a task is not the same as automating a job, and replacing a worker is not the same as replacing the judgment surrounding the work.
  \end{keyidea}
  \vspace{0.14in}
  \begin{itemize}
    \item Which parts of a job are predictable enough to automate?
    \item Which parts depend on context, accountability, trust, or tacit knowledge?
    \item What happens when organizations automate faster than they understand the failure modes?
    \item If automation fails, who absorbs the cost?
  \end{itemize}
\end{frame}

% ============================================================
\sectionframe{Part 6: Bigger Questions}
% ============================================================

\begin{frame}[shrink=4]{The Search For The Simplest Life Form}
  \begin{columns}[T,onlytextwidth]
    \begin{column}{0.46\textwidth}
      \begin{block}{Biology's Question}
        What is the smallest set of mechanisms required for something to remain alive, reproduce, and adapt?
      \end{block}
      \begin{itemize}
        \item Minimal genomes.
        \item Synthetic cells.
        \item Core mechanisms versus accumulated complexity.
      \end{itemize}
    \end{column}
    \begin{column}{0.48\textwidth}
      \begin{block}{AI Analogy}
        What is the smallest collection of mechanisms required for behavior we would call intelligent?
      \end{block}
      \begin{activity}
        Is there a ``minimum viable mind''? What capabilities would have to be included?
      \end{activity}
    \end{column}
  \end{columns}
\end{frame}

\begin{frame}[shrink=3]{Mapping The Brain}
  \begin{columns}[T,onlytextwidth]
    \begin{column}{0.47\textwidth}
      \begin{block}{Neuroscience}
        Researchers map biological brains at increasingly fine levels: regions, cell types, neurons, and connections.
      \end{block}
      \begin{itemize}
        \item The brain is extraordinarily complex.
        \item Even small nervous systems can be difficult to fully characterize.
      \end{itemize}
    \end{column}
    \begin{column}{0.47\textwidth}
      \begin{block}{Question For AI}
        Does understanding biological intelligence tell us how to engineer artificial intelligence?
      \end{block}
      \begin{itemize}
        \item Copy the mechanism?
        \item Reproduce the function another way?
        \item Learn principles without copying biology literally?
      \end{itemize}
    \end{column}
  \end{columns}
\end{frame}

\begin{frame}{Limits Of Formal Systems}
  \begin{block}{Gödel's Incompleteness Theorems}
    Roughly: sufficiently expressive formal systems contain statements whose truth cannot be established using only the system's own proof rules, assuming the system satisfies the required conditions.
  \end{block}
  \vspace{0.12in}
  \begin{keyidea}
    This is a theorem about formal mathematical systems, not a theorem that proves human minds are non-computational or that AI has a particular ceiling.
  \end{keyidea}
\end{frame}

\begin{frame}{Limits Of Formal Systems: Discussion}
  \begin{activity}
    Why are arguments about the limits of formal reasoning so tempting in discussions of intelligence? What additional assumptions would be needed to connect them to minds?
  \end{activity}
\end{frame}

% ============================================================
\sectionframe{Part 7: Course Logistics}
% ============================================================

\begin{frame}{Course Logistics And Syllabus}
  \begin{columns}[T,onlytextwidth]
    \begin{column}{0.48\textwidth}
      \begin{block}{Canvas}
        Announcements, assignments, grades, due dates, and course documents.
      \end{block}
      \begin{block}{Lectures}
        Slides are a scaffold. Class discussion, worked examples, and questions are part of the course content.
      \end{block}
    \end{column}
    \begin{column}{0.46\textwidth}
      \begin{activity}
        Open the syllabus on Canvas. We will walk through grading, late work, collaboration, AI-tool policy, labs, and major deadlines.
      \end{activity}
    \end{column}
  \end{columns}
\end{frame}

\begin{frame}[shrink=4]{About Me}
  \begin{columns}[T,onlytextwidth]
    \begin{column}{0.54\textwidth}
      \begin{itemize}
        \item Full-Time Lecturer in Computer Science and Software Engineering.
        \item Education in applied mathematics, computer science, and computer engineering.
        \item Research interests: deep reinforcement learning, heuristic search, robotics, and computer science education.
        \item Industry experience in robotics engineering and machine learning engineering.
        \item Mentoring through Beacons, OUDI, CENG SURP, research projects, and career preparation.
      \end{itemize}
    \end{column}
    \begin{column}{0.38\textwidth}
      \begin{block}{Outside Class}
        Basketball, weightlifting, video games, anime, and thinking too much about computers.
      \end{block}
      \begin{keyidea}
        I am here to help you learn the material and navigate the field.
      \end{keyidea}
    \end{column}
  \end{columns}
\end{frame}

\begin{frame}{Community}
  \begin{columns}[T,onlytextwidth]
    \begin{column}{0.42\textwidth}
      \begin{block}{Color Coded by ColorStack}
        A Cal Poly community centered on computing, belonging, peer support, professional development, and access to opportunities.
      \end{block}
    \end{column}
    \begin{column}{0.5\textwidth}
      \begin{itemize}
        \item Join the Discord.
        \item Follow \texttt{@colorcodedslo}.
        \item Watch for events, resources, and opportunities.
        \item Bring a friend.
      \end{itemize}
    \end{column}
  \end{columns}
\end{frame}

\begin{frame}{Waitlist And Roll Call}
  \begin{itemize}
    \item If you are on the waitlist, keep attending while enrollment settles.
    \item I will communicate what I can do as seats and department constraints become clear.
    \item Make sure I know you are here and that your enrollment status is correct.
  \end{itemize}
\end{frame}

\begin{frame}{Before Next Time}
  \begin{itemize}
    \item Read the syllabus.
    \item Confirm you can access Canvas and the course materials.
    \item Bring one example of a system you think behaves intelligently.
    \item Be ready to defend why you think it qualifies as AI.
  \end{itemize}
  \vspace{0.16in}
  \begin{keyidea}
    Next time: intelligent agents, environments, percepts, actions, and rational behavior.
  \end{keyidea}
\end{frame}

\begin{frame}{Next Time}
  \begin{center}
    {\Huge Intelligent Agents}
  \end{center}
  \vspace{0.15in}
  \begin{activity}
    Join me in lab.
  \end{activity}
\end{frame}

\end{document}
